\chapter{申请材料准备与文书}
\label{ch05}
\ist{材料准备}
\begin{keybox}
材料是保研里最容易掉链子的一环：成绩、排名、推荐信、政审、体检分散在教务处、院系党组织、医院和几位老师手里，任何一项卡住都可能让你错过截止时间。本章按你实际要交的东西组织——个人陈述与简历、材料清单与盖章、推荐信、政审与体检、成绩单、排名证明、证明材料、作品集。

需要先建立一个认知：\textbf{保研材料没有全国统一的清单。} 教育部层面只规定了两件全国性的事：体检依据《普通高等学校招生体检工作指导意见》（教学〔2003〕3 号，并配套教学厅〔2010〕2 号取消乙肝项目检测）；入学后 3 个月内按《普通高等学校学生管理规定》（教育部令第 41 号）复查资格，弄虚作假取消学籍。除此之外，交什么、交几份、盖哪枚章、个人陈述写多少字、推荐信要不要密封，\textbf{全部由目标院校当年简章决定}。

因此本章的做法是：凡涉及份数、字数、盖章、密封的具体要求，一律标注\textbf{学校、学院与年份}并引简章原文，做成对照表供你直接比对，而不是给你一份“保研通用清单”。看到与你报考院系不符的，以你院系当年简章为准。
\end{keybox}

\section{个人陈述}
\begin{qa}{M0258}{个人陈述一般写多少字？有没有模板}
个人陈述没有全国统一字数，各校在当年简章里各自规定，实测区间大致在 1000—3000 字。\textbf{先查目标院系当年简章，再决定写多长}，不要拿一份稿子通用。

\par\vspace{2pt}\nobreak\begin{center}\footnotesize
\begin{tabularx}{\textwidth}{>{\raggedright\arraybackslash}X>{\raggedright\arraybackslash}X}
\toprule
\textbf{学校／学院（依据年份）} & \textbf{个人陈述要求} \\ \midrule
浙江大学传媒与国际文化学院（2026） & 3000 字左右，含对申请学科专业的认识、拟定研究计划 \\
四川大学文学与新闻学院（2025） & 约 1500 字 \\
清华大学五道口金融学院（2025） & 不超过 1500 字 \\
兰州大学文学院（2026） & 不少于 1500 字 \\
山西大学电力与建筑学院（2027） & 1500 字左右 \\
合肥工业大学计算机与信息学院（2026） & 1500 字以内 \\
复旦大学物理学系（2026） & 1000 字左右 \\
北京航空航天大学仪器科学与光电工程学院（2027） & 1000 字左右，有模板（附件 7） \\
山东大学物理学院（2026） & 控制在 1000 字以内 \\
北京协和医学院群公学院（2026） & 1000 字以内，需手写签名 \\
\bottomrule
\end{tabularx}
\end{center}\vspace{2pt}\par

模板方面：北京大学力学与工程科学学院、北京大学材料科学与工程学院、北京航空航天大学仪器科学与光电工程学院、上海交通大学人文学院等都在报名系统或通知附件里提供\textbf{官方个人陈述模板}，要求下载填写后本人手写签字、不能使用电子签名。所以“有没有模板”取决于学校——有的直接给，有的只给格式不限的要求。
\nref{S0027}
\nref{S0249}
\nref{S0094}
\nref{S0004}
\nref{S0098}
\nref{S0344}
\nref{S0421}
\end{qa}

\iin{浙江大学}
\ith{个人陈述}
\ist{材料准备}
\begin{qa}{M0259}{个人陈述怎么写才不像流水账}
不像流水账的关键是\textbf{跟着简章给的内容字段走}，而不是按时间顺序罗列经历。多所学校的简章其实已经把“该写什么”写出来了：

\begin{itemize}
  \item 浙江大学传媒与国际文化学院（2026）要求含“对申请学科专业的认识、拟定研究计划”
  \item 兰州大学文学院（2026）要求含“本科期间的专业学习与科研创新情况及研究生阶段的学习计划”
  \item 北京协和医学院群公学院（2026）要求含“个人学术背景、曾经做过的研究工作、未来学习和研究计划及目标”
  \item 山东大学物理学院（2026）要求“关于曾经做过的研究工作以及未来学习和研究的计划、目标”
  \item 四川大学文学与新闻学院（2025）要求“学术背景、在所申请的专业曾经作过的研究工作、以及攻读研究生阶段的学习和研究计划”
\end{itemize}

把这几条拼起来就是通行骨架：\textbf{学术背景 → 做过什么研究（具体到问题、方法、你的贡献）→ 读研想做什么}。流水账的典型症状是只有“课程名＋社团名＋奖项名”，没有研究问题和你的动作。每个经历都回答“我解决了什么问题、用了什么方法、结果是什么、学到什么”，篇幅自然落在研究上而不是奖状上。
\nref{S0027}
\nref{S0094}
\nref{S0337}
\nref{S0385}
\nref{S0249}
\end{qa}

\iin{浙江大学}
\ith{个人陈述}
\ist{材料准备}
\begin{qa}{M0260}{个人陈述里要不要写自己为什么选这个学校}
\textbf{要写，而且这是区分度最高的一段。} 部分学校直接把“申请理由”写进个人陈述的必备字段，说明它被当作正式评分内容而非客套话：

\begin{itemize}
  \item 东北师范大学（2027）直博个人陈述明确包含“研究计划、学术背景、研究经历、\textbf{申请理由}、参与的科研项目简介及自己在其中的贡献（需要举证）”
  \item 浙江大学传媒与国际文化学院（2026）要求含“\textbf{对申请学科专业的认识}、拟定研究计划”
  \item 北京协和医学院群公学院（2026）要求含“未来学习和研究计划及目标”
\end{itemize}

写法上不要写成“贵校历史悠久、师资雄厚”。有效的版本是具体到\textbf{人、方向、平台}：你读过该院哪位老师近三年的哪项工作、和你的哪段经历接得上、你能用上学院的哪个实验室或数据。这就是“为什么是这所学校”而不是“为什么读研”。若该校简章未列申请理由字段，也不必另起一段，把这段并入“研究计划”里即可——通常两三句就够，写多了反而挤掉研究内容。
\nref{S0607}
\nref{S0027}
\nref{S0337}
\end{qa}

\iin{东北师范大学}
\ith{个人陈述}
\ist{材料准备}
\begin{qa}{M0261}{简历应该一页还是两页？要放照片吗}
\textbf{主流是一页，个别学校放宽到两页}，简章怎么写就怎么来：

\par\vspace{2pt}\nobreak\begin{center}\footnotesize
\begin{tabularx}{\textwidth}{>{\raggedright\arraybackslash}X>{\raggedright\arraybackslash}X}
\toprule
\textbf{学校／学院（依据年份）} & \textbf{简历页数要求} \\ \midrule
上海交通大学凯原法学院（2026） & 内容控制在一页 A4 纸 \\
上海交通大学上海高级金融学院（2027） & 中英文简历，正反面不超过一页 A4 纸 \\
湖南大学计算机学院（2027） & 不超过 1 页 A4 \\
中国药科大学（2027） & 限 1 页，应包含科研经历及获奖情况 \\
浙江大学建筑工程学院（2027） & A4 纸 1—2 页 \\
上海交通大学人工智能学院（2026） & 2 页内 \\
上海交通大学船舶海洋与建筑工程学院（2027） & 不超过 2 页 A4 纸，字数不限 \\
同济大学环境科学与工程学院（2027） & 面试时携带 7 份个人简历，1 页 A4 纸 \\
\bottomrule
\end{tabularx}
\end{center}\vspace{2pt}\par

照片方面：\textbf{多数简章不强制}，但北京师范大学教育学部（2026）明确要求“个人简历（含个人照片、教育背景、学习经历、获奖情况、电话号码、微信号、个人常用电子邮箱等）”。若该校写了“含个人照片”就必须放；没写则可用可不用，放了要保证与报名系统里上传的证件照同一底版。
\nref{S0240}
\nref{S0272}
\nref{S0129}
\nref{S0417}
\nref{S0025}
\nref{S0514}
\nref{S0575}
\nref{S0313}
\nref{S0151}
\end{qa}

\iin{上海交通大学}
\ith{简历}
\ist{材料准备}
\begin{qa}{M0262}{简历上没东西写怎么办，怎么填充}
先说结论：\textbf{简章只列字段，不管你有多少东西可写}，所以“没东西写”通常是把可写的东西漏掉了。先对照简章确认哪些算数：

\begin{itemize}
  \item 西安电子科技大学计算机科学与技术学院（2026）把“发表或\textbf{已录用}的、能体现自身学术水平的学术论文、出版物或原创性工作成果”与获奖证书并列作为申报材料
  \item 上海交通大学人工智能学院（2026）报名时要“详细写明作者排序、\textbf{录用状态（正在投稿／已录用）}”——说明在投论文也算
  \item 中国药科大学（2027）简历“应包含科研经历及获奖情况”
  \item 北京师范大学教育学部（2026）简历需含教育背景、学习经历、获奖情况
\end{itemize}

按这个口径盘点：课程大作业、毕设、实验室打杂、开源贡献、竞赛、课程助教、数据标注与复现实验，都属“科研经历”。真正要避免的是堆砌与专业无关的兼职和社团职务。如果确实项目少，就把一个项目写深——写清背景、你负责的模块、用了什么方法、产出是什么，比五个一句话项目更有说服力。
\nref{S0121}
\nref{S0514}
\nref{S0417}
\nref{S0151}
\end{qa}

\iin{西安电子科技大学}
\ith{简历}
\ist{材料准备}
\begin{qa}{M0263}{简历要不要写高考成绩和本科学校层次}
\textbf{简章几乎没有要求写高考成绩的，本科学校层次则通常已经在系统填报里体现了。} 从各校对简历字段的规定看：

\begin{itemize}
  \item 中国药科大学（2027）要求简历“限 1 页，应包含\textbf{科研经历及获奖情况}”
  \item 北京师范大学教育学部（2026）要求简历含“教育背景、学习经历、获奖情况、电话号码、微信号、个人常用电子邮箱”
  \item 浙江大学建筑工程学院（2027）只写“个人简历（A4 纸 1—2 页）”
\end{itemize}

这些字段清单里都没有高考成绩。本科学校层次方面，报名系统本身就会让你填学校名称，北京大学力学与工程科学学院（2026）还特别提示“系统中所填写的'所在学校'必须与成绩单上的校名一致”，上海交通大学人工智能学院（2026）要求“填写学校名称时请具体到分校”——也就是说，\textbf{学校信息由系统字段承担，不需要在简历里再强调一遍}。

因此建议：简历只放本科院校与专业、GPA／排名、科研、论文、获奖、技能。高考成绩除非该校明确要求（未见实例），不要写；主动标注“985／211／双一流”这类标签也缺少必要，评审看的是成绩单和排名证明。
\nref{S0417}
\nref{S0151}
\nref{S0025}
\nref{S0264}
\nref{S0514}
\end{qa}

\iin{中国药科大学}
\ith{简历}
\ist{材料准备}
\begin{qa}{M0275}{申请材料需要邮寄纸质版吗？截止时间按寄出还是按到达}
\textbf{两种截止口径都存在，而且“按寄达”更常见，必须逐校确认。} 从简章原文看：

\begin{itemize}
  \item 北京大学力学与工程科学学院（2026 夏令营）：纸质版“请使用 EMS，在 6 月 22 日（周一）16:30 前\textbf{显示被签收}”，推荐信需在 16 点前快递被签收
  \item 北京大学生命科学学院（2027）：“网上报名和申请材料\textbf{寄（送）达}截止时间为 2026 年 9 月 8 日 10:00，材料邮寄只接收 EMS 快递方式”
  \item 北京大学中国语言文学系（2026）：材料接收截止 2025 年 9 月 7 日 16 点，“逾期\textbf{寄送达}不予受理”
  \item 北京大学环境与能源学院（2027）：材料需在 2026 年 9 月 8 日 15:00 前送达
\end{itemize}

反过来，“以寄出时间为准（以邮戳为准）”的表述在推免材料里主要出现在\textbf{体检报告}上，例如厦门大学（2026）规定体检报告“于公示拟录取名单后 10 个工作日内提交（\textbf{以寄出时间为准}）”。

另外注意：越来越多学校已经\textbf{只要电子版、不邮寄}，浙江大学建筑工程学院（2027）明确“无需邮寄纸质材料”，北京大学材料科学与工程学院（2026）只要求把材料合成一个 PDF 发邮箱。先看清该校要不要纸质版，再决定寄什么。
\nref{S0264}
\nref{S0348}
\nref{S0095}
\nref{S0310}
\nref{S0530}
\nref{S0025}
\nref{S0380}
\end{qa}

\iin{北京大学}
\ith{邮寄}
\ist{材料准备}
\begin{qa}{M0276}{材料寄丢了怎么办，能不能补寄}
\textbf{先看该校到底要不要纸质版——这一步就能排除一半风险。} 现在大量院系已改为纯电子提交：浙江大学建筑工程学院（2027）“无需邮寄纸质材料”，北京大学材料科学与工程学院（2026）要求把所有材料按顺序整理成\textbf{一个 PDF}，命名为“意向研究所编号-2026夏令营-姓名-本科学校”，作为附件发送，并注明“勿打压缩包”。

必须寄纸质版的，简章通常写得很硬：

\begin{itemize}
  \item 北京大学力学与工程科学学院（2026）“只网上报名而不提供材料的，\textbf{视为自动放弃}”，且“因报名材料较多，不提供电话确认”
  \item 北京大学中国语言文学系（2026）“逾期寄送达不予受理”
  \item 华北电力大学（2027）“所有材料一经提交，\textbf{恕不退还}”
\end{itemize}

所以寄丢后的现实路径是：一，立刻用 email 联系院系研究生教务（简章末尾都留了邮箱和电话）说明情况，问能否补寄或先发电子版；二，如果该校明确“逾期不予受理”，补寄通常不接受，能救的只有你手上留的电子扫描件能否被接受。\textbf{提前扫描全套电子版，是唯一真正可靠的补救手段}；北京大学多所院系要求纸质材料必须与扫描件完全一致，扫描件本来就在你手上。
\nref{S0025}
\nref{S0380}
\nref{S0264}
\nref{S0095}
\nref{S0415}
\end{qa}

\iin{浙江大学}
\ith{邮寄}
\ist{材料准备}
\begin{qa}{M0279}{材料造假被查出来会有什么后果}
后果比多数人想的严重，而且是\textbf{分阶段递进}的，全国性依据明确：

\begin{itemize}
  \item 《普通高等学校招生违规行为处理暂行办法》（教育部令第 36 号）第十一条：考生“伪造、非法获得证件、成绩证明、荣誉证书等，骗取报名资格”或“在综合素质评价、相关申请材料中提供虚假材料、影响录取结果的”，\textbf{报名阶段发现取消报考资格；入学前发现取消入学资格；入学后发现取消录取资格或者学籍；毕业后发现由教育行政部门宣布学历、学位证书无效，责令收回或予以没收}；并如实记入考试诚信档案
  \item 《普通高等学校学生管理规定》（教育部令第 41 号）第十一条：学生入学后 \textbf{3 个月内}进行复查，复查中发现弄虚作假、徇私舞弊的，确定为复查不合格，\textbf{应当取消学籍}；情节严重的移交有关部门调查处理
  \item 《中华人民共和国学位法》第三十七条：以非法手段取得入学资格的，经学位评定委员会决议\textbf{不授予学位或撤销学位}
\end{itemize}

院校层面对齐执行，例如河海大学（2027）把“申请者提供的材料与事实不符，存在弄虚作假情况”列为取消录取资格的四种情形之一；中国人民大学国学院（2027）写明“无论何时，一经发现，将取消报名资格并向申请人所在高校通报”。\textbf{“入学后 3 个月内复查”意味着风险不会随录取结束而消失。}
\nref{S0041}
\nref{S0040}
\nref{S0455}
\nref{S0178}
\nref{S0006}
\end{qa}

\iin{中华人民共和国教育部}
\ith{材料风险}
\ist{材料准备}
\begin{qa}{M0280}{个人陈述找人代写靠谱吗，会不会被查重}
\textbf{简章没有针对个人陈述查重的规定，但规则环境对“代写”是明确不利的。} 需要分清三类依据：

\begin{itemize}
  \item 全国性规章中“代写”被点名的是\textbf{学位论文与研究成果}：《中华人民共和国学位法》第三十七条规定“学位论文或者实践成果被认定为存在代写、剽窃、伪造等学术不端行为”的，不授予学位或撤销学位（这与入学申请文书不是同一对象）；《高等学校预防与处理学术不端行为办法》（教育部令第 40 号）同样针对论文买卖与代写
  \item 招生环节的兜底是“材料真实性”：浙江大学传媒与国际文化学院（2026）写明“若发现申请人有弄虚作假行为，取消其申请资格并通报申请人所在本科院校”
  \item 能对得上号的技术防范出现在\textbf{推荐信}上：北京大学基础医学院（2026）明确“推荐信内容\textbf{不能雷同}”，须专家本人亲笔签字
\end{itemize}

也就是说，个人陈述代写目前主要靠面试追问与文风比对暴露，而不是查重系统。但一旦在面试中被追问细节答不上来，或与其他材料互相矛盾，就落到“材料不实”的处理轨道上。\textbf{风险不在查重，在面试。}
\nref{S0455}
\nref{S0042}
\nref{S0027}
\nref{S0301}
\end{qa}

\iin{教育部学生服务与素质发展中心（学信网/研招网）}
\ith{材料风险}
\ist{材料准备}
\begin{qa}{M0282}{申请直博要额外准备什么材料}
\textbf{直博的额外材料集中在两样：专家推荐信和研究计划／个人陈述。} 对照表如下：

\par\vspace{2pt}\nobreak\begin{center}\footnotesize
\begin{tabularx}{\textwidth}{>{\raggedright\arraybackslash}X>{\raggedright\arraybackslash}X}
\toprule
\textbf{学校（依据年份）} & \textbf{直博额外材料} \\ \midrule
东北师范大学（2027） & 申请学科或相近学科\textbf{两名正高职称}专家推荐信；个人陈述（含研究计划、学术背景、研究经历、申请理由、科研项目贡献举证） \\
华北电力大学（2027） & 至少\textbf{两名教授}（或相当专业技术职称专家）的书面推荐意见；考生个人陈述 \\
哈尔滨工业大学（2027） & 报考学科相关领域的\textbf{教授}推荐信两封 \\
华中科技大学（2027） & 《攻读博士学位研究计划》；两位\textbf{副教授及以上}专家推荐信 \\
浙江大学建筑工程学院（2027） & 至少两位\textbf{副教授及以上}专家推荐信（拟录取后 10 月底前寄纸质版） \\
四川大学建筑与环境学院（2027） & 《直接攻读博士学位研究生入学申请表》；两名\textbf{正高}职称专家推荐书，其中一份必须为报考导师出具 \\
中国人民大学财政金融学院（2027） & 2 份专家推荐书，推荐专家必须是报考领域\textbf{正高职称} \\
上海交通大学凯原法学院（2026） & 直博营另需提交\textbf{研究计划书}一份 \\
\bottomrule
\end{tabularx}
\end{center}\vspace{2pt}\par

要点：硕士阶段不少院系不要推荐信，\textbf{直博几乎一定要}；且职称门槛明显抬高（正高／教授更常见），个别学校还要求其中一封来自未来导师。
\nref{S0607}
\nref{S0415}
\nref{S0066}
\nref{S0198}
\nref{S0025}
\nref{S0070}
\nref{S0006}
\nref{S0240}
\end{qa}

\iin{东北师范大学}
\ith{直博意向}
\ist{材料准备}
\begin{qa}{M0284}{跨专业申请需要补交什么说明材料}
\textbf{跨专业申请的额外材料各校写法不同，有三种常见路径：}

\begin{itemize}
  \item \textbf{要求提交能力证明}：南京大学商学院（2026）实施细则规定“跨专业申请者应具备所申请专业大学本科阶段应达到的知识能力和水平，并须提供相关证明材料，证明材料包括\textbf{成绩单、发表学科相关论文、参与课题}等”
  \item \textbf{要求辅修／课程证明}：山东大学经济研究院（2026）与经济学院（2026）都把“\textbf{跨专业学生辅修经济学类课程证明}”列入其他材料
  \item \textbf{复试加试专业课}：中国人民大学（2027）工作办法写明“跨一级学科申请免试攻读研究生的申请人，一般应在考核时\textbf{加试一门专业课}……加试成绩不计入总分，\textbf{不及格不予录取}”；东北师范大学（2027）公费师范类办法同样规定跨专业考生须按报考院部要求加试
\end{itemize}

也有院系直接欢迎跨专业而不设附加材料，如北京大学国际法学院（2026）夏令营写明“本科专业不限，欢迎跨专业申请”。\textbf{先查目标院系是“欢迎跨专业”还是“要举证＋加试”，两类准备量差很多。}
\nref{S0282}
\nref{S0026}
\nref{S0360}
\nref{S0287}
\nref{S0609}
\nref{S0342}
\end{qa}

\iin{南京大学}
\ith{跨专业}
\ist{材料准备}
\begin{qa}{M0290}{材料里写错了信息怎么补救}
关键看\textbf{报名系统是否已截止、材料是否已被审核}，各校处理方式差别很大：

\begin{itemize}
  \item 北京大学力学与工程科学学院（2026）给了明确流程：材料\textbf{未审定前}可登录系统，进入个人信息后点“选择营地”修改志愿信息；\textbf{如果材料已经审定}，需联系咨询电话或发邮件\textbf{取消审定}后再行修改
  \item 辽宁大学（2027）：“如需修改填报信息，请再次扫描二维码登录系统，修改既有信息并点击'提交申请'即可”
  \item 冶金与生态工程学院（2026）：材料不符合要求会\textbf{退回修改}，“请在学院规定的时间内重新上传（登录系统查看），\textbf{逾期视为自愿放弃}”
  \item 西安交通大学经济与金融学院（2026）：“所有材料\textbf{一经提交不得修改}，请务必仔细检查后再提交”
  \item 中国人民大学财政金融学院（2027）：报名窗口 7 月 1 日 10:00—7 月 6 日 10:00，“\textbf{逾期系统关闭无法进行补报名}”
\end{itemize}

所以写错信息后的行动顺序是：先看系统还能不能改 → 不能改就立刻邮件联系院系教务说明错误并询问可否在审核阶段更正 → 若院系已退回材料，务必在它给的时限内重新上传。\textbf{不要沉默等待}，逾期通常不给补报机会。
\nref{S0264}
\nref{S0168}
\nref{S0270}
\nref{S0312}
\nref{S0018}

操作顺序：先截图保存当前申请状态和错误字段；再查系统是否开放修改、学院是否已审核、通知规定的材料截止时刻；若按钮不可用，尽快通过招生通知上的官方电话和邮箱报告报名号、错误内容、正确内容及证明文件，明确询问是否可退回更正。不要重复提交新账号或自行替换已盖章材料；收到书面处理答复后再操作，并保存邮件回执。
\end{qa}

\iin{北京大学}
\ith{材料风险}
\ist{材料准备}
\begin{qa}{M0298}{申请材料里的自我评价怎么写才不空}
自我评价写不空的办法是\textbf{用简章给的字段去填，而不是用形容词去堆}。多所学校把自我评价类材料的要求写得很具体：

\begin{itemize}
  \item 复旦大学物理学系（2026）：“个人自述（包括个人简历、本科期间学习科研工作简介、对未来读研和职业规划的初步设想，1000 字左右）”
  \item 华中科技大学（2027）：“个人简历或自述材料”，并附有《推免生申请个人自述模板》
  \item 中国社会科学院大学法学院（2027）把“个人简历等\textbf{自我评价材料}”与“申请攻读研究生报告（自拟）”分开列
  \item 北京协和医学院群公学院（2026）：“个人陈述（应包括个人学术背景、曾经做过的研究工作、未来学习和研究计划及目标等，1000 字以内）”
\end{itemize}

这些字段的共同点是\textbf{可被验证}。所以写“学习能力强”没有信息量，改成“大三自学了 XX 方法并用于 XX 分析，结果被 XX 采纳”就有。空洞的根因通常不是词汇问题，而是缺具体事实：先把三件做过的事写实（问题、方法、产出、你的角色），自我评价自然落到句尾，一两句收束即可。
\nref{S0344}
\nref{S0198}
\nref{S0243}
\nref{S0337}
\end{qa}

\iin{复旦大学}
\ith{个人陈述}
\ist{材料准备}
\begin{qa}{M0299}{材料里的研究计划或学习计划要不要写，写了会不会被挑刺}
\textbf{看学校：有的单独要研究计划，有的把它并进个人陈述或自述，单独的往往限定字数。} 具体分两类：

\textbf{单独作为一项材料要求的：}
\begin{itemize}
  \item 中南大学商学院（2027）：“研究生期间学习和研究计划（在系统内直接填写，\textbf{不超过 800 字}）”
  \item 四川大学化学学院（2026）：“硕士研究生期间研究计划书”
  \item 北京师范大学文理学院（2026）：“个人自述及未来研究计划（报名系统下载个人自述模板）”
  \item 华中科技大学（2027）直博生：“《攻读博士学位研究计划》”
\end{itemize}

\textbf{写在个人陈述里的：} 浙江大学传媒与国际文化学院（2026）要求个人陈述含“拟定研究计划”；北京协和医学院群公学院（2026）要求含“未来学习和研究计划及目标”；兰州大学文学院（2026）要求含“研究生阶段的学习计划”。

至于“写了会不会被挑刺”——从简章口径看，研究计划本身就是\textbf{被评分的内容}，不是可选项；面试里被追问方向细节也是常规环节。真正容易被挑的是空泛的大题目（“研究人工智能在医疗中的应用”），安全的写法是收窄到一个具体问题，并说明你已知的方法和你打算怎么入门。
\nref{S0080}
\nref{S0091}
\nref{S0149}
\nref{S0297}
\nref{S0198}
\nref{S0027}
\nref{S0337}
\end{qa}

\iin{中南大学}
\ith{研究计划}
\ist{材料准备}
\begin{qa}{M0301}{邮寄材料用顺丰还是EMS，有没有学校指定}
\textbf{有学校明确指定承运方，寄之前必须看简章那一段。} 实测到三种口径：

\begin{itemize}
  \item \textbf{只收 EMS：} 北京大学生命科学学院（2027）“材料邮寄只接收 EMS 快递方式”；北京大学力学与工程科学学院（2026）“请使用 \textbf{EMS 快递，不是挂号信}”，并要求在指定时刻前显示被签收
  \item \textbf{EMS 或顺丰二选一，且点名禁用一家：} 浙江大学能源工程学院（2026）“必须使用 EMS 或顺丰邮寄（\textbf{严禁使用京东快递}，有特殊情况无法通过 EMS 或顺丰邮寄的，请提前联系老师说明情况）”
  \item \textbf{自行选择但需保价／可追踪：} 许多院系只写“快递寄达”，不指定承运方，另有院系（如厦门大学教育研究院 2026）指定顺丰
\end{itemize}

结论：\textbf{北大多数院系要 EMS，浙大要 EMS／顺丰，个别院系指定顺丰。}“顺丰还是 EMS”没有统一答案，简章写什么就寄什么，写了“只收 EMS”却寄顺丰，风险由你自己承担。同时注意“挂号信”和“普通快递”经常被明确排除，寄件时把收件人、电话、邮编按简章原文填写，并保留单号。
\nref{S0348}
\nref{S0264}
\nref{S0542}
\nref{S0218}
\end{qa}

\iin{北京大学}
\ith{邮寄}
\ist{材料准备}
\section{材料清单}
\begin{qa}{M0248}{夏令营申请材料一般要准备哪些？有清单吗}
\textbf{没有全国统一的申请材料清单，各校差异极大。} 下表是四所学校的真实清单，可以看出共同内核和差异：

\par\vspace{2pt}\nobreak\begin{center}\footnotesize
\begin{tabularx}{\textwidth}{>{\raggedright\arraybackslash}X>{\raggedright\arraybackslash}X}
\toprule
\textbf{学校／学院（依据年份）} & \textbf{材料清单要点} \\ \midrule
北京大学力学与工程科学学院（2026 夏令营） & 申请表（手写签字、盖院系教务章）、个人陈述（系统模板）、前 5 学期成绩单（教务盖章）、英语水平证明、获奖证书复印件、论文／出版物、申请注意事项确认表、声明、专家推荐信 2 份 \\
北京大学国际法学院（2026 夏令营） & 申请表、中英文简历各 1 份、中英文个人陈述各 1 份、英语成绩报告原件、中英文成绩单、前两年半成绩排名证明、在读证明、专家推荐信 2 封、获奖证书扫描件 \\
浙江大学建筑工程学院（2027） & 个人简历（A4 1—2 页）、身份证正反面、最新专业成绩排名证明、本科成绩单（教务或院系盖章）、外语水平证明、论文与获奖证书 \\
西安电子科技大学计算机科学与技术学院（2026 预推免） & 诚信承诺书（手写签名）、个人简历、成绩单（教务盖章）、前 5 或 6 学期排名（教务盖章）、身份证及学生证复印件、外语成绩证明、其他证明材料（含已录用论文） \\
华北电力大学（2027） & 身份证／学生证复印件、推免申请表原件、政审表原件、加盖教务处公章的 1—6 学期成绩单原件、四／六级或专业外语证书复印件、科研与获奖材料复印件、诚信承诺书原件 \\
\bottomrule
\end{tabularx}
\end{center}\vspace{2pt}\par

共同内核是：\textbf{身份学籍证明 ＋ 成绩单 ＋ 排名证明 ＋ 外语证明 ＋ 个人陈述／简历 ＋ 承诺书}；差异集中在推荐信、政审表、作品与邮寄方式。
\nref{S0264}
\nref{S0342}
\nref{S0025}
\nref{S0121}
\nref{S0415}
\end{qa}

\iin{北京大学}
\ith{材料清单}
\ist{材料准备}
\begin{qa}{M0264}{申请人材料里哪些是硬性的哪些是加分项}
\textbf{多数简章会明确区分“必须提交”和“可选／如有”。} 最直白的一例是北京大学基础医学院（2026）夏令营通知：“上述材料中，\textbf{除第 6 项（其他有关能力的证明材料，如发表论文、专利、各类获奖等）由申请人依据自身实际情况提交，其余 7 项均为必须提交材料}。”

据此归纳：

\textbf{硬性项（缺失通常直接影响资格）}
\begin{itemize}
  \item 身份证、学生证或在学证明／学籍在线验证报告
  \item 加盖教务或院系公章的成绩单
  \item 成绩排名证明
  \item 外语水平证明（多数院系设 CET-6 425／500 分或同等门槛）
  \item 报名申请表／诚信承诺书（本人手写签名）
  \item 部分院系加政审表、专家推荐信
\end{itemize}

\textbf{加分项（“如有”“可选”）}
\begin{itemize}
  \item 已发表或已录用的论文、专利、出版物、原创性成果
  \item 各级获奖证书
  \item 科研项目经历、社会实践
\end{itemize}

两个需要注意的例外：一是\textbf{推荐信在部分院系是硬性的}（如北京大学多所院系要 2 封副教授及以上），在另一些院系只对直博生要求（如清华大学 2027 办法对硕士生也列为必交第 5 项）；二是本校生常被豁免若干项，清华大学（2027）明确“第 2、5、6 三项\textbf{本校应届本科毕业生不需提供}”。
\nref{S0301}
\nref{S0491}
\nref{S0121}
\nref{S0254}
\end{qa}

\iin{北京大学基础医学院}
\ith{材料清单}
\ist{材料准备}
\begin{qa}{M0267}{各种证明材料要盖什么章？学院章和教务章区别在哪}
\textbf{章的类型跟着材料走，简章通常逐项写明，不能互相替代。}

\par\vspace{2pt}\nobreak\begin{center}\footnotesize
\begin{tabularx}{\textwidth}{>{\raggedright\arraybackslash}X>{\raggedright\arraybackslash}X}
\toprule
\textbf{材料} & \textbf{常见盖章要求} \\ \midrule
成绩单 & 加盖学校\textbf{教务处公章}或\textbf{成绩单证明专用章}（北京大学环境与能源学院 2027；北京大学中国语言文学系 2026） \\
成绩排名证明 & 由所在院校\textbf{教务部门}盖章（北京大学国际法学院 2026）；也有学校接受\textbf{学院公章}（兰州大学文学院 2026 要求加盖学院公章） \\
申请表 & 所在\textbf{学校教务处或院系负责人签字盖章}；北京大学基础医学院 2026 明确“不能使用电子签名或人名章” \\
政审表／思想政治品德鉴定表 & \textbf{二级院（系）党组织公章}；浙江大学能源工程学院 2026 要求“考生所在单位政审意见”必须加盖二级院（系）党组织公章 \\
推荐信 & 专家\textbf{亲笔签名}，密封并在\textbf{封口骑缝处签字} \\
\bottomrule
\end{tabularx}
\end{center}\vspace{2pt}\par

学院章与教务章的区别在于\textbf{证明什么}：教务章证明成绩与学籍数据的权威性，学院章／党组织章证明思想政治表现和院内推荐意见。简章写“教务处公章”时，学院章通常不被接受；反之写“学院或院系公章”时两者常可通融。北京大学力学与工程科学学院（2026）还提示：申请表需“盖所在院系教务章”，而工学院本院学生可免盖——\textbf{豁免规则也逐校逐院不同。}
\nref{S0310}
\nref{S0095}
\nref{S0342}
\nref{S0094}
\nref{S0301}
\nref{S0542}
\nref{S0264}
\end{qa}

\iin{北京大学环境与能源学院}
\ith{盖章流程}
\ist{材料准备}
\begin{qa}{M0268}{盖章要跑好几个部门，有没有省事的办法}
省事的关键是\textbf{按部门而不是按学校来组织盖章路线}，因为不同学校的材料清单里，同类文件往往盖同一枚章：

\begin{itemize}
  \item 教务处／院系教务：成绩单、成绩排名证明——北京大学环境与能源学院（2027）成绩单要求“加盖学校教务处公章或成绩单证明专用章”；兰州大学文学院（2026）排名证明要求“加盖学院公章”
  \item 院系党组织：政审表／思想政治品德鉴定表——浙江大学能源工程学院（2026）要求加盖二级院（系）党组织公章
  \item 推荐专家本人：推荐信亲笔签名并密封骑缝签字——北京大学环境与能源学院（2027）
\end{itemize}

可操作的省事做法：\textbf{把同一部门的材料合并成一次跑}，成绩单和排名证明同时办、多开几份；申请表和承诺书的签字盖章尽量与成绩单同批完成。另外留一条后路——北京大学力学与工程科学学院（2026）的常见问题解答里说明“申请表暂时无法盖章的，以事实为准，可以先不盖章，之后再补盖”，合肥工业大学数学学院（2026）也注明排名证明须盖教务或院系公章；说明在\textbf{非核心材料上，部分院系允许先提交后补盖}，但成绩单这类核心材料不建议赌这一点。
\nref{S0310}
\nref{S0094}
\nref{S0542}
\nref{S0264}
\nref{S0261}
\end{qa}

\iin{北京大学环境与能源学院}
\ith{盖章流程}
\ist{材料准备}
\begin{qa}{M0277}{所有学校材料要求都不一样，有没有办法批量准备}
\textbf{能批量，但要分成“通用底稿”和“逐校替换层”两层来做。} 从多校清单看，稳定出现的共同项是：身份证／学生证、成绩单、成绩排名证明、外语水平证明、个人简历、个人陈述、诚信承诺书；差异项集中在推荐信份数与职称门槛、政审表、申请表模板、密封与邮寄方式。

\begin{itemize}
  \item 浙江大学建筑工程学院（2027）与西安电子科技大学计算机科学与技术学院（2026）的清单几乎同构：简历、身份证、排名证明、成绩单、外语证明、论文与获奖
  \item 华北电力大学（2027）多出《政审表》原件和《诚信承诺书》原件
  \item 北京大学多所院系把推荐信列为必交，并要求密封骑缝签字
\end{itemize}

具体做法：

\begin{enumerate}
  \item 建一个总目录，按“身份学籍／成绩／外语／文书／成果／推荐信／政审”分文件夹，全部存原件扫描 PDF。
  \item 成绩单、排名证明、身份证、外语成绩这些\textbf{不变项}一次办足（多开 3—5 份）。
  \item 个人陈述、简历做成主稿＋每校定制段落；申请表、承诺书、政审表按各校模板现填。
  \item 每校建一个 checklist，提交前逐项打钩，重点核对命名规则和文件大小上限。
\end{enumerate}
\nref{S0025}
\nref{S0121}
\nref{S0415}
\nref{S0301}
\nref{S0310}
\end{qa}

\iin{浙江大学}
\ith{材料清单}
\ist{材料准备}
\begin{qa}{M0278}{推荐信成绩单这些能不能一次开多份备用}
\textbf{能，而且应该一次多开——但要按用途分开准备，因为纸质版的去向不同。}

需要注意的规则：

\begin{itemize}
  \item 成绩单常被要求\textbf{密封}：北京大学环境与能源学院（2027）要求成绩单“加盖学校教务处公章或成绩单证明专用章后，装入自备信封密封，并在封口骑缝处加盖教务处公章”
  \item 推荐信必须由专家\textbf{密封并在封口骑缝处签字}（北京大学环境与能源学院 2027；清华大学 2027 办法对书面版推荐信同样要求）
  \item 纸质材料往往\textbf{不退}：华北电力大学（2027）“所有材料一经提交，恕不退还”
  \item 部分材料是\textbf{先电子后纸质}：清华大学经济管理学院（2027）要求专家在线完成推荐信后保留原件，“待后续接收通知将纸质推荐信原件寄（或送）至我院硕士项目办公室”；浙江大学建筑工程学院（2027）直博生需在 10 月底前将两份纸质专家推荐信寄至学院归档
\end{itemize}

所以建议：成绩单开 3—5 份并全部走密封流程；推荐信请专家签 2—3 份原件（很多专家愿意），密封后分别保管；其余的获奖、外语等复印件按需多印。\textbf{推荐信是唯一无法自己补开的材料，多签几份的边际成本最低。}
\nref{S0310}
\nref{S0254}
\nref{S0025}
\nref{S0415}
\end{qa}

\iin{北京大学环境与能源学院}
\ith{材料清单}
\ist{材料准备}
\begin{qa}{M0283}{申请专硕和学硕的材料有区别吗}
\textbf{学硕与专硕的差别主要体现在论文／成果形式上，真正的材料大分化在“直博 vs 硕士”这条线上。}

专硕／学硕的差异，西南财经大学国际商学院（2027）写得最清楚：报考\textbf{学术型硕士}的学生提交的论文应为学术论文；报考\textbf{专业型硕士}的学生“论文的形式不限于学术论文，也可采用\textbf{案例分析、调研报告、研究计划}等形式”；而报考直博项目的考生“需要另外提交 5 份\textbf{研究计划书}”。

直博带来的额外材料在多数学校是一致的：至少两名专家推荐信（职称门槛普遍更高），加研究计划或攻读博士学位研究计划。例如：

\begin{itemize}
  \item 中国人民大学财政金融学院（2027）直博项目要求 2 份专家推荐书，推荐专家必须取得\textbf{正高职称}；同时报考专硕和直博的，若进入直博拟推荐范围，\textbf{不再获得专硕项目考核资格}
  \item 四川大学建筑与环境学院（2027）直博需另交《直接攻读博士学位研究生入学申请表》与两名正高职称专家推荐书，其中一份必须为报考导师出具
  \item 上海交通大学人工智能学院（2026）要求报名时慎重填写报考类别（直博或直硕），\textbf{报名后无极特殊情况不得变更}
\end{itemize}

也就是说，硕士阶段学硕／专硕的材料清单通常一致，差异在评审偏好；报直博则要额外准备推荐信和研究计划。
\nref{S0320}
\nref{S0018}
\nref{S0070}
\nref{S0514}
\end{qa}

\iin{西南财经大学国际商学院}
\ith{材料清单}
\ist{材料准备}
\begin{qa}{M0287}{电子签名或电子章的材料外校认不认}
\textbf{分两种情况：学校官方系统出具的电子版通常认，个人签名的电子化替代通常不认。}

\textbf{认可的}：上海交通大学船舶海洋与建筑工程学院（2027）“对于加盖公章的材料需为扫描件，或\textbf{优先使用官方部门开具的电子版文件}（如学校教务部门开具的电子版成绩单、成绩排名证明、电子版英语成绩单等）”；北京大学力学与工程科学学院（2026）也把“学校系统申请的电子成绩单，打印出来加盖教务章”列为合格成绩单之一。

\textbf{不认可的}：北京大学基础医学院（2026）对申请表和个人陈述都写明“须由申请人\textbf{本人手写签字}……\textbf{不能使用电子签名或人名章}”；北京大学力学与工程科学学院（2026）要求“本人手写签字，\textbf{不能使用电子签名}”；北京大学中国语言文学系（2026）对推荐信要求“亲笔签名（电子签名或复印签名等形式的签名均无效）”。

实操坑：北京大学力工学院提示“\textbf{电子成绩单不能直接合并进 PDF，建议打印出来以后扫描，否则电子章将无法显示}”。
\nref{S0575}
\nref{S0264}
\nref{S0301}
\nref{S0095}
\end{qa}

\iin{上海交通大学}
\ith{盖章流程}
\ist{材料准备}
\begin{qa}{M0288}{申请材料里的照片有什么要求}
照片要求分两层：\textbf{推免服务系统的统一要求}和\textbf{各校报名系统的附加要求}。

教育部推免服务系统（学信网／研招网，2024 年发布的推免生信息准备说明）的要求是全国性的：

\begin{itemize}
  \item 本人\textbf{近三个月内}正面、免冠、\textbf{无妆}、彩色头像电子证件照（蓝色或白色背景）
  \item 仅支持 \textbf{jpg 或 jpeg} 格式，建议大小不超过 10M，\textbf{宽高比例 3:4}
  \item 正脸头像，人像水平居中；人脸水平转动角、倾斜角、俯仰角在 ±10 度内；眼睛位置距照片上边沿为图像高度的 30\%—50\%
  \item 头部和肩部占整张照片比例\textbf{不小于 2/3}；脸部无遮挡，露出五官
  \item \textbf{不得化妆，不得佩戴眼镜、隐形眼镜、美瞳}
\end{itemize}

各校附加要求更严，例如：

\begin{itemize}
  \item 北京师范大学（2026）推免申请系统：彩色证件照，\textbf{高 200 像素、宽 150 像素}，jpg，\textbf{30k 以内}
  \item 北京大学基础医学院（2026）：单独上传 1 张\textbf{1 寸证件照}（JPG，小于 300kb），命名为“姓名+申报专业名称”
  \item 北京协和医学院群公学院（2026）：个人证件照片\textbf{限 100K 内}，jpg，以“XX 大学-XX 专业-XX（姓名）”命名
\end{itemize}

\textbf{系统照片一旦定稿通常影响后续所有材料，务必先按系统要求拍，再按各校尺寸压一遍。}
\nref{S0453}
\nref{S0423}
\nref{S0301}
\nref{S0337}
\end{qa}

\iin{教育部学生服务与素质发展中心（学信网/研招网）}
\ith{材料清单}
\ist{材料准备}
\begin{qa}{M0291}{申请时用的邮箱有什么讲究？QQ邮箱可以吗}
\textbf{邮箱本身极少被简章限制，真正决定成败的是“邮件主题和附件命名”。}

从各校简章原文看，要求几乎都指向命名：

\begin{itemize}
  \item 华中科技大学网络空间安全学院（2026）：“Email 至邮箱……（邮件主题和文件名命名为'2026 夏令营-申请人姓名'），\textbf{未按要求命名可能导致邮件丢失或被过滤}”
  \item 北京林业大学外语学院（2026）：“邮件主题及压缩包均命名为'2026 年夏令营-大学名称-姓名'”
  \item 北京大学信息管理系（2026）：“邮件主题和 PDF 文件附件命名均为'出版硕士/图书情报硕士+本科学校+姓名+报名材料'”
  \item 贵州大学计算机科学与技术学院（2027）：“包命名格式：考生就读学校-考生姓名-报考专业”
\end{itemize}

也就是说，\textbf{规范命名比邮箱品牌更重要}。至于 QQ 邮箱能不能用：简章未见任何禁用条款，实际使用中也能正常投递；但要注意两点——一，部分高校邮件网关会拦截邮件，所以按简章要求的主题格式命名能显著降低被过滤概率；二，长期沟通（联系导师、接收录取通知）建议用\textbf{稳定、日常查看的邮箱}，并在简历和报名系统里写同一个。
\nref{S0003}
\nref{S0023}
\nref{S0035}
\nref{S0002}
\end{qa}

\iin{华中科技大学}
\ith{材料清单}
\ist{材料准备}
\begin{qa}{M0292}{材料打包命名有没有规范，会不会影响印象}
\textbf{会影响印象，而且部分学校把命名写成了硬性要求，写错可能直接被过滤。}

实测到的命名格式（都来自简章原文）：

\par\vspace{2pt}\nobreak\begin{center}\footnotesize
\begin{tabularx}{\textwidth}{>{\raggedright\arraybackslash}X>{\raggedright\arraybackslash}X}
\toprule
\textbf{学校（依据年份）} & \textbf{要求的命名规则} \\ \midrule
北京大学材料科学与工程学院（2026） & PDF 命名为“意向研究所编号-2026夏令营-姓名-本科学校”，作为附件发送，\textbf{勿打压缩包} \\
西安电子科技大学计算机科学与技术学院（2026） & “计科院+预推免+本人姓名+本人所在高校” \\
北京大学信息管理系（2026） & “出版硕士/图书情报硕士+本科学校+姓名+报名材料” \\
河南大学外语学院（2026） & “2026夏令营+本科专业+姓名” \\
北京林业大学外语学院（2026） & “2026年夏令营-大学名称-姓名” \\
贵州大学计算机科学与技术学院（2027） & “考生就读学校-考生姓名-报考专业” \\
\bottomrule
\end{tabularx}
\end{center}\vspace{2pt}\par

可以看出\textbf{通行结构是“学校/专业-姓名-学校名”这一类组合}，用短横线或加号连接，不用空格和特殊符号。另外两条通用规则：一，\textbf{按要求合并成 1 个 PDF}，不要把文件散着发；二，\textbf{除非简章说可以压缩，否则不要打成 zip}（北京大学材料学院明确“勿打压缩包”）。命名不规范虽不必然扣分，但会让审材料的人多花时间，也会提高被邮件系统过滤的概率。
\nref{S0380}
\nref{S0121}
\nref{S0035}
\nref{S0399}
\nref{S0023}
\nref{S0002}
\end{qa}

\iin{北京大学}
\ith{材料清单}
\ist{材料准备}
\begin{qa}{M0293}{申请同一学校多个学院要重复交材料吗}
\textbf{多数学校限制的不是“重复交材料”，而是“同一单位只能报一个志愿”。} 两者的区别很大——能不能同时申请同一所学校的两个学院，取决于该校系统怎么设。

实测的限制条款：

\begin{itemize}
  \item 中国人民大学财政金融学院（2027）：“申请人在本学院直博专业\textbf{只能选一个}，同时申请两个以上专业者，\textbf{材料不予受理}”；另一则通知同样写“申请人在本学院只能申请一个专硕项目，同时申请两个以上者，材料不予受理”
  \item 中国人民大学国学院（2027）：“每位申请人\textbf{限报一个类别及一个专业方向}，同时申请我院两个及两个以上类别及专业方向者将视为无效报名”
  \item 北京大学信息科学技术学院（2026）：“申请者\textbf{只能申请院内一个志愿}”
  \item 浙江大学（2027）：“每个申请人\textbf{只能填报一个志愿}”
  \item 北京语言大学语言科学与资源学院（2026）：“每名夏令营申请人\textbf{限报一个招生单位、一个专业}”
  \item 中南大学商学院（2027）：“申请人仅限报名以上一个专业（研究方向）”
\end{itemize}

所以同一学校内跨学院同时申请，很多学校在系统层面就不允许；即便允许，材料也要\textbf{按各院要求分别准备、分别提交}，不存在“交一次两院通用”。查清该校报名系统的志愿限制，比准备重复材料更要紧。
\nref{S0018}
\nref{S0019}
\nref{S0006}
\nref{S0114}
\nref{S0032}
\nref{S0021}
\nref{S0080}
\end{qa}

\iin{中国人民大学}
\ith{材料清单}
\ist{材料准备}
\begin{qa}{M0294}{保研材料需不需要提前扫描存档？怎么整理}
\textbf{要，而且从准备材料的第一天就要开始扫。} 三个理由都来自简章原文：

\begin{itemize}
  \item \textbf{纸质版须与扫描件完全一致}：北京大学力学与工程科学学院（2026）写明“扫描件和开营提交的纸版资料原文件必须完全一致，否则取消夏令营资格”；北京大学信息科学技术学院等多所院系也要求“纸质报名材料必须与扫描件完全一致”
  \item \textbf{纸质件押后提交}：北京大学多个院系统一表述为“入营营员还需提交纸质报名材料和专家推荐信 2 封，须在规定时间内快递寄达学院，具体信息将另行通知，\textbf{请勿提前邮寄}”——扫描件先行、纸质件押后，两套必须对应
  \item \textbf{电子版是唯一补救凭证}：材料寄丢、系统故障需要补交时，手上有没有清晰 PDF 决定能否救回来
\end{itemize}

整理建议：按“身份学籍／成绩排名／外语／文书／成果证书／推荐信”六个子目录归档，统一 \textbf{300dpi 彩色扫描}，每份单独存 PDF 再另存一份合并版，同步云盘。
\nref{S0264}
\nref{S0114}
\nref{S0143}
\nref{S0213}
\nref{S0575}
\end{qa}

\iin{北京大学}
\ith{材料清单}
\ist{材料准备}
\section{推荐信}
\begin{qa}{M0249}{推荐信一般要几封？找谁写最合适}
\textbf{份数各校自定，2 封最常见；“找谁”则由简章规定职称与相关性。} 对照表：

\par\vspace{2pt}\nobreak\begin{center}\footnotesize
\begin{tabularx}{\textwidth}{>{\raggedright\arraybackslash}X>{\raggedright\arraybackslash}X}
\toprule
\textbf{学校（依据年份）} & \textbf{推荐信要求} \\ \midrule
北京大学多所院系（2026／2027） & 专家推荐信 \textbf{2 封}，2 位与申请学科相关的\textbf{副教授（含）以上}或相当专业技术职称专家分别推荐 \\
清华大学（2027）／清华大学经济管理学院（2027） & \textbf{2 封}，两位与申请学科相关的\textbf{副教授（含）以上}职称专家在线填写 \\
北京大学人工智能研究院（2026） & \textbf{3 封}，相关学科专业领域内副教授（含）以上 \\
合肥工业大学数学学院（2026） & \textbf{1 封} \\
合肥工业大学计算机与信息学院（2026） & \textbf{1 封}，须是\textbf{具有硕导资格}的教师 \\
华东师范大学计算机科学与技术学院（2026） & \textbf{1 封}，副教授及以上职称专家推荐 \\
北京师范大学（2026，2027） & 2 封，\textbf{仅限申请博士}（硕士不要求） \\
中国人民大学财政金融学院（2027，直博） & 2 份专家推荐书，专家必须是报考领域\textbf{正高职称} \\
上海交通大学凯原法学院（2026） & 2 封副教授（含）以上教师推荐信，\textbf{至少一封来自所在院系} \\
\bottomrule
\end{tabularx}
\end{center}\vspace{2pt}\par

选人优先级：\textbf{熟悉你 + 职称达标 + 与申请方向相关}。三样里“熟悉你”最容易被牺牲，但它决定推荐信有没有细节。
\nref{S0342}
\nref{S0254}
\nref{S0350}
\nref{S0261}
\nref{S0098}
\nref{S0108}
\nref{S0423}
\nref{S0006}
\nref{S0240}
\end{qa}

\iin{北京大学}
\ith{推荐信}
\ist{材料准备}
\begin{qa}{M0250}{推荐信必须找教授或院长吗，讲师行不行}
\textbf{普遍底线是“副教授（含）以上”，讲师一般不行；但个别简章留了口子。}

主流表述：

\begin{itemize}
  \item 北京大学多所院系：“2 位与申请攻读学位学科相关专家【\textbf{副教授职称（含）或相当专业技术职称以上}】分别签名推荐”
  \item 清华大学（2027）：“两位与申请学科相关的\textbf{副教授（含）以上}职称专家”
  \item 上海交通大学凯原法学院（2026）：“两封\textbf{副教授（含）以上}职称教师的推荐信”
  \item 南开大学马克思主义学院（2027）：“两位\textbf{高级职称（副教授/教授）}专家亲笔签名推荐信”
\end{itemize}

例外与变体：

\begin{itemize}
  \item 合肥工业大学计算机与信息学院（2026）写的是“\textbf{具有硕导资格}教师推荐信 1 封”——这种情况下讲师只要具备硕导资格就可以，限制的是\textbf{导师资格}而不是职称
  \item 中国人民大学财政金融学院（2027）直博项目要求专家必须是报考领域的\textbf{正高职称}，门槛比副教授更高
  \item 哈尔滨工业大学（2027）直博要求“报考学科相关领域的\textbf{教授}（或相当专业技术职称专家）”
\end{itemize}

所以判断标准是：\textbf{先看该校写的是“职称”还是“资格”}。写职称（副教授及以上）就找副教授起步；写硕导资格就可以找有带研资格的讲师。院长、系主任的头衔本身不构成加分，简章看的是职称和相关性。
\nref{S0342}
\nref{S0491}
\nref{S0240}
\nref{S0133}
\nref{S0098}
\nref{S0006}
\nref{S0066}
\end{qa}

\iin{北京大学}
\ith{推荐信}
\ist{材料准备}
\begin{qa}{M0251}{推荐信是自己写让老师签字吗，这样算不算造假}
\textbf{简章的要求是“由专家填写／由推荐专家密封签字”，自己写好让老师签字与这一要求相悖；但现实中这种做法相当常见，风险主要在面试核查和内容雷同上。}

简章原文的措辞值得注意：

\begin{itemize}
  \item 清华大学（2027）：“申请人在线邀请两位……专家\textbf{在线填写}推荐信。推荐专家按申请服务系统要求，在规定的时间内\textbf{填写、上传}专家推荐信”，并要求专家在系统完成后协助保存原件，后续寄送
  \item 北京大学基础医学院（2026）：推荐信“\textbf{须由专家本人亲笔签字}，推荐信内容\textbf{不能雷同}”
  \item 北京大学计算机学院（2026）：“推荐信正文内容可电脑录入，也可以手写，\textbf{必须专家亲笔签名}，不能电子签名”
\end{itemize}

看出去两层意思：简章允许正文由电脑录入，但\textbf{签名必须是专家手写}，且专家要为内容负责。自己起草、老师审阅后签字，如果老师确实读过并认可，与“内容不能雷同”的要求并不必然冲突；但如果多封推荐信内容雷同、或面试中问到推荐信里的细节你答不上来，就会出问题。\textbf{最稳的做法是提供素材提纲（项目、表现、评价角度）让老师自己写，或至少让老师逐句改过。}
\nref{S0491}
\nref{S0301}
\nref{S0114}
\nref{S0254}
\end{qa}

\iin{清华大学}
\ith{推荐信}
\ist{材料准备}
\begin{qa}{M0252}{推荐信要不要密封？现在很多是线上提交还要密封吗}
\textbf{看提交方式：线下寄纸质版的几乎都要密封骑缝签字，纯线上提交的不需要密封。}

\textbf{要密封的（纸质推荐信）：}

\begin{itemize}
  \item 北京大学环境与能源学院（2027）：“推荐信须\textbf{密封并在封口骑缝处签字}”
  \item 北京大学中国语言文学系（2026）：推荐信“须亲笔签名（电子签名或复印签名等形式的签名均无效）、\textbf{密封并在封口骑缝处签字}”
  \item 清华大学（2027）：“书面版\textbf{密封（在密封信封的封口骑缝处签字）}交申请人寄至所申请院系的研究生教学办公室”
  \item 清华大学经济管理学院（2027）：系统在线填写后，需将“纸质推荐信原件（推荐信须由推荐专家\textbf{密封并在封口骑缝处签字}）”寄至项目办公室
  \item 北京大学力学与工程科学学院（2026）：纸质版材料仅需寄送 2 封专家推荐信，用 EMS 寄出
\end{itemize}

\textbf{不需要密封的（专家在线上传）：} 清华大学与北京大学多数院系的报名系统里，推荐信由专家在\textbf{系统内在线填写、上传}，此时密封环节转移到后续纸质归档阶段；北京大学基础医学院（2026）甚至明确“\textbf{专家推荐信不用提交电子版}”，只在纸质材料里交原件。

一句话判断：\textbf{该信最终是以“信封”形式寄出去的，就要密封＋骑缝签字；只由专家在系统里提交的，系统本身就是密封。}
\nref{S0310}
\nref{S0095}
\nref{S0491}
\nref{S0254}
\nref{S0264}
\nref{S0301}
\end{qa}

\iin{北京大学环境与能源学院}
\ith{推荐信}
\ist{材料准备}
\begin{qa}{M0253}{推荐信模板都长什么样？有没有可以参考的框架}
\textbf{推荐信框架的“评分维度”其实写在简章里}，照着填比套网上模板更稳。

东北师范大学（2027）推免办法对直博推荐信的内容规定得最细：“推荐信应对申请人\textbf{思想品德、道德修养、科学作风、治学态度、外语水平、知识结构、科研能力、工作成果、培养潜力及综合素质}作出简要评价”——这十个维度基本可以当作通用的内容清单。

格式与自由度方面：

\begin{itemize}
  \item 上海交通大学凯原法学院（2026）：“两封副教授（含）以上职称教师的推荐信（\textbf{至少一封来自所在院系，无格式要求}）”
  \item 北京大学多所院系：提供\textbf{官方推荐信模板}，需从报名系统下载，专家填写后密封
  \item 北京师范大学（2026）：提供《申请北京师范大学夏令营专家推荐信》模板，正文“可由专家手写或打字，但\textbf{签字应为专家手写亲笔签名}”
\end{itemize}

所以框架建议是：开头一段说明推荐人与你的关系（课程、科研、指导时长）→ 中段按 2—3 个维度各举一个具体事例（科研能力、学习态度、协作能力）→ 结尾给比较性评价与推荐强度。\textbf{没有模板的学校，就用上面的维度自拟，别去抄网上的范本。}
\nref{S0607}
\nref{S0240}
\nref{S0301}
\nref{S0149}
\end{qa}

\iin{东北师范大学}
\ith{推荐信}
\ist{材料准备}
\begin{qa}{M0254}{老师答应写推荐信但一直拖，快截止了怎么办}
\textbf{先分清是“线上填写”还是“纸质寄送”，因为这两种的截止点不一样，补救手段也不同。}

简章对流程的规定：

\begin{itemize}
  \item 清华大学（2027）：“申请人在线邀请两位……专家\textbf{在线填写}推荐信。推荐专家按申请服务系统要求，\textbf{在规定的时间内}填写、上传专家推荐信”，之后还要请专家保存原件、后续寄送书面版
  \item 清华大学经济管理学院（2027）：同样是在线邀请，纸质原件“\textbf{待后续接收通知}”再寄
  \item 北京大学力学与工程科学学院（2026）：纸质版只寄 2 封专家推荐信，需在指定时刻前被签收，并特别注明“\textbf{2 封推荐信如遇特殊情况可分开邮寄}”
\end{itemize}

可操作的三条：

\begin{enumerate}
  \item \textbf{把系统里专家邮箱填对并提前告知老师}，很多系统会直接给专家发填写链接，老师漏看邮件是最常见的拖延原因。
  \item 距截止 5—7 天老师还没动，就换策略：当面找或电话，别只发微信；同时准备好备选推荐人。
  \item 若纸质件来不及同时寄，参考北大做法\textbf{分开邮寄}，先寄已经签好的那封并邮件告知教务，问清能否后补另一封。\textbf{主动联系院系教务说明情况，比沉默等到截止要有效得多。}
\end{enumerate}
\nref{S0491}
\nref{S0254}
\nref{S0264}
\end{qa}

\iin{清华大学}
\ith{推荐信}
\ist{材料准备}
\begin{qa}{M0255}{推荐信里能写缺点吗，写缺点会不会减分}
\textbf{简章没有要求推荐信写缺点，但也没有禁止；从简章对推荐信功能的定位看，客观评价比一味夸奖更可信。}

看两处原文：

\begin{itemize}
  \item 东北师范大学（2027）要求推荐信对申请人“思想品德、道德修养、科学作风、治学态度、外语水平、知识结构、科研能力、工作成果、\textbf{培养潜力}及综合素质作出\textbf{简要评价}”——“评价”是中性词，允许有保留
  \item 北京大学基础医学院（2026）明确“推荐信内容\textbf{不能雷同}”——用意正是防止千篇一律的溢美之词
\end{itemize}

从评审角度，一封全是“该生学习刻苦、成绩优异”的信没有任何区分度，反而会让人怀疑推荐人是否真的了解你。\textbf{可以写、也值得写的“不足”是带成长性的}：比如“该生早期在文献阅读上偏窄，但在完成 X 项目后主动补齐了 Y 方向的基础”，或者“性格偏内向，但在组会上逐渐能主动汇报”。这类表述既真实，又展示了自我修正能力。

要避免的是\textbf{方向性否定}：与申请专业直接冲突的短板（如“不适合做研究”“数学基础薄弱”）不要出现在推荐信里，简章的评分维度里没有一项是“缺点”，写了只会减分。
\nref{S0607}
\nref{S0301}
\end{qa}

\iin{东北师范大学}
\ith{推荐信}
\ist{材料准备}
\begin{qa}{M0256}{换过导师，推荐信该找谁写}
\textbf{原则是：推荐人要与“申请学科方向”相关，并尽量有一封来自你现在的院系。} 换过导师不构成障碍，但要处理好两件事——推荐人的选择，和相关性证明。

简章的相关规定：

\begin{itemize}
  \item 上海交通大学凯原法学院（2026）：“两封副教授（含）以上职称教师的推荐信，\textbf{至少一封来自所在院系}”
  \item 清华大学（2027）：在线邀请“两位与\textbf{申请学科相关}的副教授（含）以上职称专家”
  \item 西南财经大学国际商学院（2027）：直博生推荐信“其中 \textbf{1 封须为被报考博导的推荐信}”
  \item 中国人民大学财政金融学院（2027）：直博推荐专家必须是报考领域\textbf{正高职称}
\end{itemize}

据此可分三种情形处理：

\begin{itemize}
  \item \textbf{原导师仍在相关方向}：请他写第二封，另一封找现任院系里了解你的老师（满足“至少一封来自所在院系”）
  \item \textbf{原导师已不在相关方向或不便联系}：优先找带过你科研的现任老师 + 一门核心课老师
  \item \textbf{两段科研经历分属不同方向}：可以让两位老师各写一段，侧重点不同，反而比两封重复的信更有信息量
\end{itemize}

换导师本身不需要在推荐信里解释，如果需要说明，放在个人陈述而非推荐信里。
\nref{S0240}
\nref{S0491}
\nref{S0320}
\nref{S0006}
\end{qa}

\iin{上海交通大学}
\ith{推荐信}
\ist{材料准备}
\begin{qa}{M0257}{实习单位领导能写推荐信吗}
\textbf{多数简章限定的是“职称”，实习单位领导通常不具备该职称，因此一般不能作为专家推荐信提交——但要看简章用的是“职称”还是“专业技术职称/行业专家”这类表述。}

典型表述：

\begin{itemize}
  \item 北京大学国际法学院（2026）：需要 2 位“\textbf{副教授（含）以上或具有相当专业技术职称者}分别推荐”
  \item 清华大学（2027）：“两位与申请学科相关的\textbf{副教授（含）以上职称专家}”
  \item 北京大学基础医学院（2026）：“相关专业\textbf{副教授及以上或具有相当专业技术职称}的专家推荐”
  \item 上海交通大学凯原法学院（2026）：“两封\textbf{副教授（含）以上职称教师}的推荐信”
\end{itemize}

这些措辞的主体都是高校或科研机构的教研人员，实习单位的项目经理、部门主管通常不符合。“相当专业技术职称”一般指副研究员、高级工程师等\textbf{有职称序列的岗位}，如果你的实习单位领导恰好具备高级工程师等职称，且工作内容与申请方向相关，理论上可以沟通，但风险在于院系审核时可能不认。

实操建议：\textbf{实习经历写进简历和个人陈述，推荐信留给学术推荐人。} 如果实习单位与申请方向高度相关（如申请金融硕、实习在券商研究部），可以向院系教务单独询问是否接受行业推荐信，得到明确答复再提交。
\nref{S0342}
\nref{S0491}
\nref{S0301}
\nref{S0240}
\end{qa}

\iin{北京大学}
\ith{推荐信}
\ist{材料准备}
\begin{qa}{M0297}{推荐信提交后还能修改吗}
\textbf{看提交方式，也看系统是否已锁定；多数学校的答案是“一般不能改，个别情况可退回重传”。}

不同情形的原文依据：

\begin{itemize}
  \item \textbf{系统提交后锁定：} 西安交通大学经济与金融学院（2026）明确“所有材料\textbf{一经提交不得修改}，请务必仔细检查后再提交”
  \item \textbf{审核不通过可退回重传：} 冶金与生态工程学院（2026）“如材料不符合要求将\textbf{退回修改}，请在学院规定的时间内重新上传（登录系统查看），\textbf{逾期视为自愿放弃}”
  \item \textbf{材料未审定前可改：} 北京大学力学与工程科学学院（2026）说明材料未审定前可登录系统修改志愿信息，已审定则需联系咨询电话或发邮件\textbf{取消审定}后再改
  \item \textbf{由专家本人操作：} 清华大学（2027）推荐信由系统在线邀请专家填写上传，申请人无法代改；清华大学经济管理学院（2027）还要求纸质原件内容“需与系统提交的电子版保持一致，以便后续存档”
\end{itemize}

所以推荐信提交后想改，路径只有两条：一是系统允许退回（多数仅在审核不通过时发生）；二是直接联系院系研究生教务说明情况，请求让推荐人在系统中重新提交。\textbf{不要试图自己用第二个邮箱另开一封顶替}，专家邮箱与身份是对应的，会被系统或人工核查发现。
\nref{S0312}
\nref{S0270}
\nref{S0264}
\nref{S0491}
\nref{S0254}
\end{qa}

\iin{西安交通大学}
\ith{推荐信}
\ist{材料准备}
\section{体检}
\begin{qa}{M0269}{政审表什么时候要，内容是什么}
\textbf{两个时间点：夏令营／预推免阶段有的要，复试或拟录取后也有的要；内容是思想政治表现，不是体检。} 政审表在简章里常叫“思想政治品德鉴定表”“现实表现复审表”。

实测要求：

\begin{itemize}
  \item 华北电力大学（2027）：材料清单第（3）项为“《政审表》\textbf{原件}”，在\textbf{复试报到及资格审查}时提交
  \item 中国科学院沈阳自动化研究所（2026 夏令营）：申请材料里含“《研究生政审表》”，与推荐信、成绩单一同按顺序装订
  \item 中国科学技术大学天文与空间科学学院（紫金山天文台）（2027）：材料含“考生\textbf{思想政治品德鉴定表}或考生情况政审表（签字、盖章）”
  \item 浙江大学能源工程学院（2026）：政审表在\textbf{拟录取后}由本人档案所在高校二级学院（系）填写，并需单位负责人签字及单位公章
  \item 中山大学（2026）：拟录取硕士研究生在线下载“\textbf{现实表现复审表}”，按调档函要求寄回
\end{itemize}

内容一般包括：政治态度、思想表现、道德品质、遵纪守法、诚实守信等，由你档案所在单位（本科院系党组织或学生处）填写并盖章。\textbf{多数学校把它放在拟录取之后}，但夏令营和预推免阶段的院系也可能提前要，所以提前找辅导员沟通是必要的。
\nref{S0415}
\nref{S0056}
\nref{S0294}
\nref{S0542}
\nref{S0164}
\nref{S0219}
\end{qa}

\iin{华北电力大学}
\ith{政审}
\ist{材料准备}
\begin{qa}{M0270}{政审表谁签字盖章？辅导员还是学院党委}
\textbf{签字盖章的主体是“你档案所在单位”，具体到本科学校就是院系党组织（或学生处），不是辅导员个人。} 各校表述略有差异但指向一致：

\begin{itemize}
  \item 浙江大学能源工程学院（2026）：“政审表需填写完整，并含有所在单位\textbf{负责人签字及单位公章}，方为有效”，并明确“'考生所在单位政审意见'必须加盖\textbf{二级院（系）党组织公章}”
  \item 中国科学院信息工程研究所（2026 夏令营）：“政审表（附件 2，手填、打印均可），需\textbf{所在院系党总支或学生处签字盖章}”
  \item 西北工业大学（2026）：“《西北工业大学研究生招生复试政审表》……由考生本人\textbf{档案所在学校（或单位）人事、政工部门}审核，填写意见并盖章”
  \item 安徽工业大学（2026 年博士招生简章）：“拟录取名单确定后，学院应向考生所在单位函调人事档案和本人现实表现等材料”“向考生所在单位函调人事档案和本人现实表现等材料（须加盖考生人事档案所在单位\textbf{人事或政工部门公章}）”
  \item 浙江大学能源工程学院还要求“政审表应当由政审单位寄送，且需\textbf{密封并在封口处加盖单位公章}”
\end{itemize}

也就是说：辅导员可以协助办理，但落到公章上通常是\textbf{院系党组织章}（如“XX学院党委”）或学校学生工作部门的章；\textbf{团委章、学生会章、个人签名一般不被接受}。填写内容必须由政审单位填写，“报考院系”“鉴定时段”等项目也要按该校填表说明来填，浙江大学能源工程学院就明确规定填写不符合填表说明的不予接受。
\nref{S0542}
\nref{S0219}
\nref{S0506}
\nref{S0161}
\end{qa}

\iin{浙江大学}
\ith{政审}
\ist{材料准备}
\begin{qa}{M0271}{体检是什么时候做？保研体检和考研体检一样吗}
\textbf{推免体检通常在拟录取后进行，不与考研体检同步；有效期普遍按“三个月内”控制。}

时间点上的规定：

\begin{itemize}
  \item 中山大学（2027）：“所有拟录取的推免生，须于\textbf{拟录取后}在各招生单位指定的医院或自行选择符合要求的医院进行体检，体检要求及时间由招生单位另行通知”
  \item 浙江大学能源工程学院（2026）：“体检时间距寄送时间应\textbf{不超过三个月}”
  \item 厦门大学（2026）：“体检报告于\textbf{公示拟录取名单后 10 个工作日内}提交（以寄出时间为准），体检报告日期 \textbf{3 个月内有效}”
  \item 福建农林大学（2026）：“体检报告有效期以复试时间为准，\textbf{三个月内有效}”
  \item 宁波大学材料科学与化学工程学院（2027）：“考生应自行前往二级甲等（含）以上医院进行体检，并将体检报告原件（须有体检结论）于\textbf{复试成绩公布后一周内}提交给学院”
  \item 华中科技大学（2027）：“\textbf{正式报名阶段}还需提供二级甲等以上医院或各高校校医院的体检证明”
\end{itemize}

医院层级普遍要求\textbf{二级甲等及以上}，也有学校接受“各高校校医院”。体检项目参考华中科技大学列出的口径：内科、外科、胸透、五官科、肝功能两项（ALT 谷丙转氨酶、AST 谷草转氨酶）。

\textbf{结论：别在夏令营阶段就急着体检}——太早做的报告到拟录取时已经过期；等拟录取或院系通知后再去，按二甲以上医院做，留好原件。
\nref{S0167}
\nref{S0542}
\nref{S0530}
\nref{S0413}
\nref{S0153}
\nref{S0198}
\end{qa}

\iin{中山大学}
\ith{体检}
\ist{材料准备}
\begin{qa}{M0272}{体检有指标不合格会不会影响录取}
\textbf{体检不合格会被取消录取资格，但“不合格”的门槛由全国性文件界定，范围比想象中窄。} 依据《普通高等学校招生体检工作指导意见》（教学〔2003〕3 号）。

该文件规定：\textbf{“除患有传染性疾病、精神性疾病、血液病、心脏病、高血压等无法完成学业的疾病及学习不能自理的考生，高等学校可以不予录取外，对患有其他疾病的考生，只要不影响专业学习和其他学生，录取时一般应不受限制。”} 文件把受限情形分为三类——学校可不予录取的、录取受限的专业、以及“\textbf{不宜就读}”的指导性建议；第三类只是报考参考，文件明确“学校不得以此为依据，拒绝录取达到相关要求的考生”。

院校层面：厦门大学（2026）把“体检不合格者”列入不予录取情形；福建农林大学（2026）规定体检不合格者或无故未在规定时间提交体检报告者取消拟录取资格。复查依据是《普通高等学校学生管理规定》（教育部令第 41 号）第十一条，入学后 3 个月内复查“身心健康状况是否符合报考专业或者专业类别体检要求”。

\textbf{要点：影响录取的是“能否完成学业”以及专业本身的额外身体条件限制。}
\nref{S0039}
\nref{S0530}
\nref{S0413}
\nref{S0040}
\end{qa}

\iin{中华人民共和国教育部}
\ith{体检}
\ist{材料准备}
\begin{qa}{M0273}{乙肝携带或色盲等情况会被卡吗}
\textbf{乙肝检测已被全国性文件取消，不作为录取依据；色盲、色弱则要看具体专业是否受限。}

\textbf{乙肝}：依据《教育部办公厅 卫生部办公厅关于普通高等学校招生学生入学身体检查取消乙肝项目检测有关问题的通知》（\textbf{教学厅〔2010〕2 号}），入学体检取消乙肝项目检测。中山大学（2027）、厦门大学（2026）等多校在体检条款里同时援引教学〔2003〕3 号与教学厅〔2010〕2 号执行。

\textbf{色盲／色弱}：教学〔2003〕3 号按色觉异常程度列了受限专业，其中\textbf{色觉异常 II 度（色盲）}不能录取的专业包括美术学、绘画、艺术设计、摄影、动画、博物馆学、应用物理学、天文学、地理科学、应用气象学、材料物理、矿物加工工程、资源勘探工程、冶金工程、无机非金属材料工程、交通运输、油气储运工程等。

\textbf{结论：乙肝携带本身不构成录取障碍；色盲色弱要对照报考专业逐条核对教学〔2003〕3 号的专业限制名单，并向目标院系确认。}
\nref{S0167}
\nref{S0530}
\nref{S0638}
\nref{S0039}
\end{qa}

\iin{中山大学}
\ith{体检}
\ist{材料准备}
\begin{qa}{M0274}{无犯罪记录证明保研需要吗}
\textbf{在本手册采集的 639 份高校官方来源（含 2026、2027 年推免办法与院系细则）里，没有检索到任何一所高校要求推免申请人提交《无犯罪记录证明》。} 各校要求的材料是另一套东西：

\begin{itemize}
  \item 华北电力大学（2027）材料清单第（3）项是“《政审表》原件”
  \item 中国科学技术大学天文与空间科学学院（紫金山天文台）（2027）要求“考生\textbf{思想政治品德鉴定表}或考生情况政审表（签字、盖章）”
  \item 浙江大学能源工程学院（2026）用《浙江大学拟录取研究生\textbf{思想政治素质和品德鉴定表}》（简称政审表），需所在高校二级学院（系）填写并加盖党组织公章
  \item 西北工业大学（2026）要求《研究生招生复试\textbf{政审表}》由档案所在学校（或单位）人事、政工部门审核盖章
  \item 中山大学（2026）拟录取后在线下载“\textbf{现实表现复审表}”
\end{itemize}

也就是说，\textbf{功能上对应的是“思想政治素质和品德考核”，通过政审表／品德鉴定表完成，而不是公安机关出具的无犯罪记录证明}。另外，本科学校在\textbf{推免资格遴选}阶段可能把“无违法犯罪记录”写成遴选条件之一——例如成都大学商学院推免细则把“品行表现优良，无违法犯罪记录和尚未解除的纪律处分”列为基本条件，但这是校内遴选资格的门槛，不是要你交一张证明。

由于所查公开资料无法穷尽所有学校，\textbf{建议以你报考院系当年简章的材料清单为准；如果清单里没写，就不必主动准备。}
\end{qa}

\ith{政审}
\ist{材料准备}
\section{成绩单}
\begin{qa}{M0243}{保研成绩单要去哪里开？教务处自助打印机开的够用吗}
\textbf{能在自助打印机开，但“够不够用”有明确标准——北大给出了三种合格形式。} 北京大学力学与工程科学学院（2026 夏令营）常见问题解答原文：

\begin{itemize}
  \item 纸质版打印成绩单，\textbf{加盖学校或者院系教务（章）}的
  \item \textbf{自助打印机}打印的直接带学校或者院系教务章（或成绩证明专用章），且\textbf{使用特殊纸张}打印出来的（自助打印机的纸张一般是防伪纸）
  \item 学校系统申请的\textbf{电子成绩单}，打印出来一致后，\textbf{加盖学校或者院系教务（章）}的
\end{itemize}

并特别提示：“\textbf{电子成绩单不能直接合并进 PDF，建议打印出来以后扫描，否则电子章将无法显示。}”

其他学校方向一致：北京师范大学（2026）要求“成绩单（须加盖学校教务部门或院系公章，扫描为 PDF 文件，大小 1M 以内）”；河南大学历史文化学院（2026）要求教务部门出具的加盖公章的成绩单与总成绩排名证明各 1 份（由学院或教务部门打印盖章）。

\textbf{结论：自助打印机开出的成绩单只要带教务章／成绩证明专用章且用防伪纸，就是合格形式；纯电子成绩单直接拼 PDF 会因电子章不显示而失效。}
\nref{S0264}
\nref{S0423}
\nref{S0250}
\end{qa}

\iin{北京大学}
\ith{成绩单}
\ist{材料准备}
\begin{qa}{M0244}{成绩单要不要英文版？什么时候需要}
\textbf{中文成绩单是刚需，英文成绩单只在特定场合需要——主要是全英文授课项目、部分经管类项目，以及中外合作办学／海外交换背景的申请。} 实测要求：

\begin{itemize}
  \item 北京大学国际法学院（2026 夏令营，全英文授课）：材料第 5 项为“本科阶段\textbf{中文成绩单原件}（由教务部门盖章）和\textbf{英文成绩单}（\textbf{如院方不出具英文成绩单，可自行翻译，成绩需和盖章的中文成绩单保持一致}）”，可见需中英双份，且接受自行翻译
  \item 清华大学经济管理学院（2027）金融硕士等项目：要求提交“本科\textbf{中、英文学业成绩单}原件扫描件，攻读第二学位或参加过海外交换的学生也请提供相应成绩单原件扫描件”
\end{itemize}

另一条容易踩的时间线来自学校自己的规定：北京大学教务部（2027）说明中提到“已获得推免资格但未落实接收单位的，推免资格失效；\textbf{推免接收工作结束后，方可办理英文成绩单}”——也就是说，本科学校开具英文成绩单的窗口期可能和推免时间线重叠甚至在后，不要在推免截止前临时去要。

\textbf{建议：如果你报的是全英文项目或有海外交换经历，提前问本科教务处英文成绩单的开具周期；纯中文项目一般不要求英文版。}
\nref{S0342}
\nref{S0254}
\nref{S0553}
\end{qa}

\iin{北京大学}
\ith{成绩单}
\ist{材料准备}
\begin{qa}{M0245}{成绩单上显示挂科重修记录有影响吗，能不能申请不显示}
\textbf{简章不允许“申请不显示”——成绩单必须真实完整。}

\textbf{真实性层面}：北京大学力学与工程科学学院（2026）在成绩单一项后注明“学校或者院系教务盖章成绩单，\textbf{如有不实将取消入营资格}”；北京大学材料科学与工程学院（2026）要求“包含所修\textbf{全部课程}的正式成绩单原件”。申请不显示在操作上等于伪造成绩单，落入《普通高等学校招生违规行为处理暂行办法》（教育部令第 36 号）第十一条的处理轨道。

\textbf{记录本身的影响}：一是最后一学年，中山大学（2027）把“\textbf{本科最后一学年课程成绩出现不及格记录}”列为取消录取资格情形，浙江大学（2027）表述为“本科最后一学年计入毕业要求学分的相关课程成绩出现不及格记录”；二是入门门槛，北京协和医学院群公学院（2026）要求“\textbf{无任何不及格、补考记录}”。

\textbf{早前学期的挂科若已重修通过，成绩单会同时显示原始与重修记录，这是常态。真正要守住的是最后一学年不再出现新的不及格。}
\nref{S0264}
\nref{S0380}
\nref{S0167}
\nref{S0032}
\nref{S0337}\par\textbf{本版实操补充：} 不能自行删改或要求出具不真实版本。先向教务处确认成绩单能否同时展示首次成绩与重修成绩、是否有官方的推免用格式；提交时按招生单位要求提供完整盖章件。解释时简洁说明补修通过时间、后续相关课程表现和改进方式。若细则把“曾有不及格”列为硬门槛，应在投递前向招生单位确认，不要用隐去记录的材料冒险。
\end{qa}

\iin{北京大学}
\ith{成绩单}
\ist{材料准备}
\begin{qa}{M0296}{成绩单上有些课是P或NP，外校看得懂吗}
\textbf{外校看得懂 P／NP 与等级制，但你需要主动说明换算口径。}

多数高校对非百分制成绩有明文折算规则。成都大学商学院推免细则直接给出对照表：五级制“优秀／良好／中等／及格／不及格”对应百分制 \textbf{95／85／75／65／55}，并说明“如果某课程的原始成绩采用的是五级制，则按照以下对应规则转换为百分制分数后再纳入上述计算公式计算”。

另一类做法是把解释权交给申请人。北京大学力学与工程科学学院（2026）的常见问题解答写明：“成绩排名是专业课排名还是所有修过课程的成绩排名，请与本科院系联系确认，\textbf{如果仍无法确认，可以自行填写，并在其它说明中注明}”——学校允许你\textbf{附说明}，但不允许留下歧义。院校评判口径也提示了风险：北京协和医学院群公学院（2026）要求“平均学分绩点对应成绩 \textbf{85 分及以上}”。

\textbf{建议：在成绩单之外附一页说明}，写清 P/NP 是“通过／不通过”还是“未计入 GPA”、等级制与百分制的换算规则、以及你是否把 P/NP 课程计入加权。
\nref{S0079}
\nref{S0264}
\nref{S0337}
\end{qa}

\iin{成都大学商学院}
\ith{成绩单}
\ist{材料准备}
\begin{qa}{M0300}{海外本科的成绩单和评分体系怎么转换成国内认可的形式}
\textbf{没有全国统一的“转换”规则，通行做法是“原版成绩单＋官方认证＋自行翻译说明”。}

\begin{itemize}
  \item \textbf{认证文件是硬要求}：教育部《全国硕士研究生招生工作管理规定》明确，考生在录取当年入学前须取得“教育部留学服务中心出具的《国（境）外学历学位认证书》”，未取得者\textbf{录取资格无效}
  \item \textbf{可自行翻译}：北京大学国际法学院（2026）要求英文成绩单“如院方不出具英文成绩单，\textbf{可自行翻译}，成绩需和盖章的中文成绩单保持一致”
  \item \textbf{海外交换经历也要交}：清华大学经济管理学院（2027）要求提交中、英文学业成绩单原件扫描件，“攻读第二学位或参加过海外交换的学生也请提供相应成绩单原件扫描件”
  \item \textbf{优先官方电子版}：上海交通大学船舶海洋与建筑工程学院（2027）“优先使用\textbf{官方部门开具的}电子版文件（如学校教务部门开具的电子版成绩单、成绩排名证明）”
\end{itemize}

做法：留服认证提前办；保留学校官方成绩单原件扫描；另附一页换算说明，写清 GPA 制度、满分标准和排名口径。
\nref{S0049}
\nref{S0051}
\nref{S0342}
\nref{S0254}
\nref{S0575}
\end{qa}

\iin{中华人民共和国教育部}
\ith{成绩单}
\ist{材料准备}
\section{排名证明}
\begin{qa}{M0246}{排名证明必须写专业排名吗？学院只肯开年级排名怎么办}
\textbf{简章通常写“成绩排名证明”，但排名口径要你自己在与本科学校核对后标注清楚——有的学校同时要专业排名和综合排名。}

\begin{itemize}
  \item 北京大学基础医学院（2026）写得最细：“成绩排名证明(\textbf{含专业成绩排名和综合成绩排名}。因排名系统不同，请依照所在本科院校排名的规定填写，由所在学校教务部门盖章或院系出具证明；请与所在本科院校联系确认，成绩排名是\textbf{专业课排名还是所有课程的成绩排名}，请在申请表上备注说明)”
  \item 北京大学力学与工程科学学院（2026）：“系统中的综合排名，根据自身情况如实填写，\textbf{如所在学校尚无综合排名，填写后另附说明}”
  \item 合肥工业大学数学学院（2026）：“专业成绩排名：\textbf{无硬性要求}，但需提供成绩排名证明；或有突出成果、高水平论文或特殊专长者（需在申请表的显著位置注明）”
  \item 河南大学外语学院（2026）要求“本科阶段成绩单与\textbf{总评成绩排名证明}（需要加盖就读院系公章）”
  \item 山东大学经济研究院（2026）：“前五学期的成绩单原件与总评成绩排名证明（均需单位盖章），\textbf{如果成绩排名暂时没有，九月份补寄亦可}”
\end{itemize}

\textbf{学院只肯开年级排名时的处理}：照实提交年级排名，并在申请表备注或另附说明中写清“我院仅按年级口径出具排名，专业方向为 XX，专业人数 XX”。评审看的是口径是否清楚、有无掩饰，而不是排名口径本身。
\nref{S0301}
\nref{S0264}
\nref{S0261}
\nref{S0399}
\nref{S0026}
\end{qa}

\iin{北京大学基础医学院}
\ith{排名证明}
\ist{材料准备}
\begin{qa}{M0247}{排名证明上的排名和成绩单对不上，外校会怀疑吗}
\textbf{会怀疑，而且有的学校直接把“不一致”写成了无效报名。} 最直接的一条是中南大学商学院（2027）的规定：“所在院(系)或学校教务处盖章的成绩单原件的扫描件（\textbf{系统填报的排名须与教务盖章排名一致，不一致者视为无效报名}）”。

其他相关要求：

\begin{itemize}
  \item 上海交通大学凯原法学院（2026）要求排名“加盖所在学校教务处或所在学院教务办的公章，\textbf{同一院校申请人应采用同一排名规则}”——这条同时约束了你和你的同学
  \item 北京大学力学与工程科学学院（2026）说明：“如果本科院校学校只有'综合排名'或者只有'成绩排名'，那另外一项填写的时候，\textbf{可以一样的排名，并在其它说明中注明}（比如'学校只提供成绩排名，综合排名填写与成绩排名一致'）”；还提示“如果仍无法确认，可以自行填写，并在其它说明中注明……\textbf{总之，以院系盖章数据为准}”
  \item 北京大学基础医学院（2026）要求排名证明含“专业成绩排名和综合成绩排名”，并“请在申请表上备注说明”口径
\end{itemize}

\textbf{所以对不上的正解不是改数字，而是解释口径}：系统里填的排名、盖章证明上的排名、成绩单算出来的排名，三者不一致时，说明“哪个是专业口径、哪个是年级口径、哪个是综合测评口径”，必要时把三种都写出来。\textbf{偷偷对齐数字而被查出前后矛盾，性质就变成不实材料了。}
\nref{S0080}
\nref{S0240}
\nref{S0264}
\nref{S0301}
\end{qa}

\iin{中南大学}
\ith{排名证明}
\ist{材料准备}
\begin{qa}{M0289}{材料提交系统崩溃或超时了怎么办}
\textbf{多数简章对逾期不给补救，所以策略是“不把提交压到最后一小时”。} 简章里的措辞普遍很硬：

\begin{itemize}
  \item 中国人民大学财政金融学院（2027）：报名时间为 7 月 1 日 10:00—7 月 6 日 10:00，“\textbf{逾期系统关闭无法进行补报名}”
  \item 中国人民大学财政金融学院（2027）：“通过材料初审的考生务必查看……在规定的时间内下载复试通知书，\textbf{逾期系统将自动关闭下载功能}”
  \item 北京协和医学院群公学院（2026）：“报名系统开放时间内完成信息填报及材料上传，\textbf{逾期系统自动关闭，不予补报}，请务必在规定时间内完成线上报名”
  \item 西安交通大学经济与金融学院（2026）：“所有材料\textbf{一经提交不得修改}，请务必仔细检查后再提交”
  \item 冶金与生态工程学院（2026）：材料有问题会退回修改，“请在学院规定的时间内重新上传（登录系统查看），\textbf{逾期视为自愿放弃}”
\end{itemize}

可操作的四条：

\begin{enumerate}
  \item \textbf{提前 48 小时完成提交}，不要卡最后一天；系统在截止前几小时最容易拥堵。
  \item 提交前把材料全部\textbf{转成 PDF 存在本地}，不依赖云端。
  \item 遇到系统故障，第一时间\textbf{打简章上的咨询电话}并同时发邮件留痕，附上截图和你的报名号。
  \item 部分学校有辅助系统（如华中科技大学 2027 的预报名系统、中国人民大学的推免生信息采集系统），\textbf{学校系统和教育部推免服务系统是两套，都要报}。
\end{enumerate}
\nref{S0018}
\nref{S0337}
\nref{S0312}
\nref{S0270}
\nref{S0198}
\end{qa}

\iin{中国人民大学}
\ith{提交事故}
\ist{材料准备}
\begin{qa}{M0295}{学院不给开排名证明怎么办}
\textbf{分两种情况：学校允不允许后补，以及能否用替代文件顶替。}

\textbf{允许后补的实例}：山东大学经济研究院（2026）明确“前五学期的成绩单原件与总评成绩排名证明（均需单位盖章），\textbf{如果成绩排名暂时没有，九月份补寄亦可}”；北京大学力学与工程科学学院（2026）常见问题解答也说明“申请表暂时无法盖章的，以事实为准，可以先不盖章，之后再补盖”。

\textbf{必须提供的}：合肥工业大学数学学院（2026）虽写“专业成绩排名：无硬性要求，但\textbf{需提供成绩排名证明}”；北京师范大学心理学部（2026）要求“本科成绩单扫描件、\textbf{专业排名证明（须加盖学校教务处或院系公章）}”。可以不要高排名，但不能不交证明。

\textbf{学院不给开时的三条路：} 一，换部门——学院不开通常可找学校教务处或学生服务大厅自助打印系统出排名证明；二，换口径——只开年级排名的照实提交并在备注说明口径（北京大学基础医学院 2026 要求“在申请表上备注说明”）；三，书面说明——由学院出具无法提供排名证明的说明并盖章，附成绩单与自算均分，向报考院系说明情况。\textbf{主动沟通好过材料缺失被退回。}
\nref{S0026}
\nref{S0264}
\nref{S0261}
\nref{S0297}
\nref{S0301}
\end{qa}

\iin{山东大学}
\ith{排名证明}
\ist{材料准备}
\section{证明材料}
\begin{qa}{M0265}{获奖证书需要盖章复印件吗？电子版怎么处理}
\textbf{多数简章只要复印件或扫描件，不要求盖章；重点是“简历／申请表里提到的每一项都要有对应扫描件”。}

实测要求：

\begin{itemize}
  \item 北京大学环境与能源学院（2027）：“获奖证书\textbf{复印件} 1 套”
  \item 北京大学国际法学院（2026）：“\textbf{简历中提及的获奖证书}扫描件各 1 份”——注意这句，意味着简历里写了什么就要附什么，写了不附反而扣印象
  \item 北京大学基础医学院（2026）：把“其他有关能力的证明材料(如：发表论文、专利，各类获奖等)”列为可由申请人自行决定提交的项，其余为必交
  \item 西安电子科技大学计算机科学与技术学院（2026）：“其他相关证明材料（如获奖证书，发表或已录用的……学术论文、出版物或原创性工作成果）”
  \item 合工大数学学院（2026）材料清单第 6 项即为“获奖证书”
\end{itemize}

也有个别要求\textbf{核验原件}，例如华南师范大学国际商学院（2026）：“其他有关材料复印件（如大学英语四、六级考试成绩单，其他获奖证书……），\textbf{需验原件}”。

\textbf{电子版处理建议}：按“简历／申请表出现的顺序”编号扫描，合并成一个 PDF，多页的证书只扫首页或有姓名和奖项名的那一页；文件命名统一为“序号-奖项名-年份”。原件在入营报到时通常要带，扫描件与原件必须一致。
\nref{S0310}
\nref{S0342}
\nref{S0301}
\nref{S0121}
\nref{S0261}
\nref{S0009}
\end{qa}

\iin{北京大学环境与能源学院}
\ith{证明材料}
\ist{材料准备}
\begin{qa}{M0266}{论文没发表，提交录用通知或投稿记录可以吗}
\textbf{多数学校接受——但要把“录用状态”写清楚，不能含糊成“已发表”。}

原文依据：

\begin{itemize}
  \item 上海交通大学人工智能学院（2026）在报名说明里直接要求：“论文请详细写明\textbf{作者排序、录用状态（正在投稿／已录用）}”——说明在投和录用都被接受，前提是标明状态
  \item 西安电子科技大学计算机科学与技术学院（2026）把“发表或\textbf{已录用}的、能体现自身学术水平的学术论文、出版物或原创性工作成果”列入其他证明材料
  \item 北京协和医学院群公学院（2026）：“发表文章（\textbf{可选，限第一、二作者}，只需文章首页并标出作者）；原创性工作成果（可选）”
  \item 北京大学环境与能源学院（2027）：“申请人还可提交体现自身学术水平的代表性学术论文、出版物或\textbf{原创性工作成果}”
\end{itemize}

\textbf{可提交的材料形态}包括：录用通知（附期刊名称、卷期或 DOI）、投稿系统状态截图、预印本（如 arXiv）、论文首页。署名方面，多数院系在评审时会关注\textbf{作者排序}，第一、二作者权重最高，部分院系干脆限定“限第一、二作者”。

\textbf{两条红线：} 一，不要写“已发表”而实际只是“已投稿”——这属于材料不实；二，未发表作品若用了他人成果或数据，同样落入学术不端范畴。写清状态、附上可核查的凭证，录用通知完全可以是加分项。
\nref{S0514}
\nref{S0121}
\nref{S0337}
\nref{S0310}
\end{qa}

\iin{上海交通大学}
\ith{证明材料}
\ist{材料准备}
\begin{qa}{M0285}{学校要求提供课程大纲或培养方案，去哪找}
\textbf{在本手册核对的 639 份来源中，没有出现要求推免申请人提交本科“课程大纲”的简章}（所查资料中的“培养方案”多指研究生培养方案或强基转段方案，与本科课程大纲不是一回事）。实际要求的是两类更轻的材料：

\begin{itemize}
  \item \textbf{跨专业课程证明}：山东大学经济研究院（2026）与经济学院（2026）把“\textbf{跨专业学生辅修经济学类课程证明}”列入其他材料，要求提供相应证明文件
  \item \textbf{已修课程成绩说明}：上海理工大学新闻传播学专业夏令营（2026）给出务实替代做法——“在校成绩单或者\textbf{已修课程成绩说明}电子文档（考虑到暑假期间，成绩单的打印、盖章存在诸多不便，请在成绩单末加注对所提供成绩真实性承诺，并由本人签字说明）”
\end{itemize}

若确实被要求提供课程内容说明（多见于国外院校申请或需要评估先修课的院系），来源有本科院系教务处的课程简介库、学院官网“课程设置与大纲”栏目、或任课教师开具的课程说明。\textbf{就国内推免而言，成绩单＋辅修证明通常就够了，不必主动准备课程大纲。}
\nref{S0026}
\nref{S0360}
\nref{S0081}
\end{qa}

\iin{山东大学}
\ith{证明材料}
\ist{材料准备}
\begin{qa}{M0286}{排名证明和成绩单需要公证吗}
\textbf{未见统一规定，各校要求不同——在本手册核对的 639 份高校官方来源（含 2026、2027 年推免办法与院系细则）中，没有任何一份文件出现“公证”字样，也没有要求排名证明或成绩单办理公证的条款。} 各校对这两类材料的实际要求是“盖章”，而不是“公证”：

\begin{itemize}
  \item 成绩单：北京大学环境与能源学院（2027）要求“加盖学校教务处公章或成绩单证明专用章”
  \item 排名证明：北京大学国际法学院（2026）要求“由所在院校教务部门盖章”；河南大学外语学院（2026）要求“加盖就读院系公章”
  \item 电子版：上海交通大学船舶海洋与建筑工程学院（2027）“优先使用\textbf{官方部门开具的}电子版文件（如学校教务部门开具的电子版成绩单、成绩排名证明）”
  \item 英文材料：北京大学国际法学院（2026）对英文成绩单的处理是“如院方不出具英文成绩单，\textbf{可自行翻译}，成绩需和盖章的中文成绩单保持一致”
\end{itemize}

\textbf{所以对国内推免而言，成绩单和排名证明的权威性来自学校教务部门的公章，不需要公证。} 需要公证的场景通常是出国留学申请、境外学历认证等，与国内推免是两套流程。本结论属阴性发现，不能替代对目标院系当年简章的逐条核对；如果报考的是中外合作办学项目或有特殊规定的专项计划，请以该项目当年的材料清单为准。
\end{qa}

\ith{证明材料}
\ist{材料准备}
\section{作品集}
\begin{qa}{M0281}{保研文书要不要做作品集？艺术类怎么做}
\textbf{未见统一规定，各校要求不同——在本手册核对的 639 份来源（含 2026、2027 年推免办法与院系细则，覆盖艺术类、建筑类、设计类院系）中，没有任何一份文件出现“作品集”这一要求。} 与之功能相近的是\textbf{“代表作”}，且主要集中在人文社科类专业：

\begin{itemize}
  \item 北京大学历史学系（2027）招生指南写明：“本系录取研究生的标准……还注重对考生学术背景材料的审查，\textbf{特别重视能表明考生研究能力的代表作品}。请报考考生提供能证明自身学术研究能力的代表成果”
  \item 北京师范大学教育学部教育学原理相关专业夏令营（2026）：材料第 4 项为“\textbf{代表作}（已发表或未发表的学术论文、读书报告、学术随笔等均可）”
  \item 北京师范大学教育学部（2026）：“本科期间本人撰写的学术论文或能证明写作能力的\textbf{文字作品}（文稿）一篇”
  \item 厦门大学教育研究院（2026）：“如申请者在报名资料中提供了公开发表的论文，请同时将其中一篇\textbf{代表作}电子版（word 版）以'姓名+2026年夏令营申请'为邮件主题发送”
\end{itemize}

\textbf{结论：所查公开资料中未见“作品集”要求，无法给出艺术类的作品集规格（页数、格式、提交方式）。} 建议你直接查目标院系当年简章——艺术与设计类院系的要求通常写在简章附件的材料清单或复试细则里；若简章未列，可在报名前邮件询问教务。
\end{qa}

\ith{作品集}
\ist{材料准备}

\begin{qa}{M0302}{多所院系要求不同材料文件名，如何避免上传错件？}
按目标院系给出的原始文件名格式优先命名，建立“学校\_学院\_项目\_材料类型\_姓名”的本地分类目录，但只在该院系允许的字段中填写姓名或编号。逐校做一张材料清单，把成绩单、排名证明、个人陈述和证书扫描件分别核对；上传前打开文件确认内容、页序、方向与清晰度，再对照申请页面逐项勾选。不同单位要求不一致时不要用一套文件覆盖全部申请，也不要把含身份证号码的文件名用于公开网盘。

再做一次提交前核验：按“学校—学院—项目—材料类型—版本日期”建立本地文件夹，目录仅自用；文件名严格服从该校模板，不把其他学校要求的命名规则带过去。上传后逐项回看申请系统中的材料栏目，确认没有把排名证明放进成绩单栏、没有误传上一版文件，并记录每个申请的提交时间和回执。含证件号码的扫描件只通过学校指定系统提交，不放在公开网盘或共享链接中。
\end{qa}

\section{5.9 多校材料版本管理}
\begin{qa}{U202610210501}{成绩单更新后，不同院系收到的新旧版本不一致怎么办？}
先核对每个申请批次要求的成绩截止时间和材料版本：有的单位要求截至前五学期，有的要求最新成绩单，不能只凭文件日期判断。建立投递记录，标注每所学校提交的成绩单版本、覆盖学期、开具日期、上传时间及是否允许补交；新成绩影响排名或满足门槛时，立即按招生单位公开渠道询问是否需要补传，并说明新旧版本差异。不要静默替换已提交文件或在不同栏目重复上传造成歧义；保存原始文件和学校答复，后续面试按实际提交版本解释。
\end{qa}
\begin{qa}{U202610210502}{线上申请显示材料已上传，怎样确认招生单位收到的是完整文件？}
上传完成提示只证明平台接受了文件，不一定代表内容清晰、栏目匹配或审核通过。提交前按学校清单逐项核对文件名、格式、大小限制、页数、签章和扫描清晰度；提交后回到申请记录，检查文件是否仍可预览、状态是否为“已提交/待审核”，并保存带时间的页面截图或系统回执。若系统只显示文件名、不提供预览，保留本地原件和上传记录，向通知指定联系人确认是否需要邮件补交。发现错误时先询问更正入口和截止时间，不要重复创建多个申请账号。\nref{S0761}
\end{qa}
